\documentclass[11pt]{article}

\usepackage[margin=0.82in,headheight=22pt]{geometry}
\usepackage{amsmath,amssymb,amsthm,bm}
\usepackage{fontspec}
\setmainfont{TeX Gyre Pagella}
\setsansfont{TeX Gyre Heros}
\setmonofont{Courier New}[Scale=0.86]
\usepackage{microtype}
\usepackage{xcolor}
\usepackage{booktabs,longtable,array,colortbl}
\usepackage{enumitem}
\usepackage{listings}
\usepackage{fancyhdr}
\usepackage{titlesec}
\usepackage{hyperref}
\usepackage{bookmark}

\definecolor{navy}{RGB}{24,55,95}
\definecolor{teal}{RGB}{20,112,104}
\definecolor{codebg}{RGB}{246,248,250}
\definecolor{rulegray}{RGB}{210,216,222}
\definecolor{codecomment}{RGB}{80,100,90}

\hypersetup{
  colorlinks=true,
  linkcolor=navy,
  urlcolor=teal,
  pdftitle={STAT452 60-Minute Written Test: Complete Solutions},
  pdfauthor={STAT452 Study Guide}
}

\pagestyle{fancy}
\fancyhf{}
\lhead{\sffamily\small STAT452 60-Minute Written Test}
\rhead{\sffamily\small Mathematical Solutions and R Verification}
\cfoot{\sffamily\thepage}
\renewcommand{\headrulewidth}{0.4pt}
\renewcommand{\headrule}{\hbox to\headwidth{\color{navy!45}\leaders\hrule height \headrulewidth\hfill}}

\titleformat{\section}{\sffamily\LARGE\bfseries\color{navy}}{\thesection}{0.65em}{}
\titleformat{\subsection}{\sffamily\Large\bfseries\color{teal}}{\thesubsection}{0.65em}{}
\titleformat{\subsubsection}{\sffamily\large\bfseries\color{navy!85}}{\thesubsubsection}{0.65em}{}
\titlespacing*{\section}{0pt}{2em}{0.75em}
\titlespacing*{\subsection}{0pt}{1.45em}{0.55em}
\titlespacing*{\subsubsection}{0pt}{1.15em}{0.4em}

\setlength{\parindent}{0pt}
\setlength{\parskip}{0.52em}
\setlength{\emergencystretch}{1.5em}
\setlist{leftmargin=1.55em,itemsep=0.25em,topsep=0.35em}
\renewcommand{\arraystretch}{1.18}

\lstset{
  language=R,
  basicstyle=\ttfamily\small,
  keywordstyle=\color{navy}\bfseries,
  commentstyle=\color{codecomment}\itshape,
  stringstyle=\color{teal!80!black},
  backgroundcolor=\color{codebg},
  frame=single,
  rulecolor=\color{rulegray},
  breaklines=true,
  columns=fullflexible,
  keepspaces=true,
  showstringspaces=false,
  xleftmargin=0.45em,
  xrightmargin=0.45em,
  aboveskip=0.6em,
  belowskip=0.7em
}

\newcommand{\E}{\operatorname{E}}
\newcommand{\Var}{\operatorname{Var}}
\newcommand{\Cov}{\operatorname{Cov}}
\newcommand{\logit}{\operatorname{logit}}
\newcommand{\SSE}{\operatorname{SSE}}
\newcommand{\SST}{\operatorname{SST}}
\newcommand{\RSS}{\operatorname{RSS}}
\newcommand{\Rtwo}{R^2}

\begin{document}
\hypersetup{pageanchor=false}

\begin{titlepage}
  \centering
  \vspace*{1.0in}
  {\sffamily\color{navy}\Huge\bfseries STAT452 60-Minute Written Test\par}
  \vspace{0.25in}
  {\sffamily\color{teal}\LARGE Complete Mathematical Solutions\par}
  \vspace{0.16in}
  {\sffamily\Large with R Verification\par}
  \vspace{0.5in}
  {\color{navy!45}\rule{0.78\textwidth}{1.1pt}\par}
  \vspace{0.5in}
  {\Large Applied Statistics for Engineers and Scientists II\par}
  \vfill
  \begin{minipage}{0.82\textwidth}
    \colorbox{navy!7}{\parbox{0.94\linewidth}{
      \textbf{Structure.} Each problem is solved first with the mathematical
      notation used in the course Notes, followed by R code that verifies the
      calculation. The ridge result in Problem 4 is proved using derivatives
      and centered sums.
    }}
  \end{minipage}
  \vfill
  {\sffamily\small Based on \texttt{MidTerm\_STAT452\_23TT2.pdf} and Notes Chapters 1--4\par}
\end{titlepage}

\hypersetup{pageanchor=true}
\pagenumbering{roman}
\tableofcontents
\clearpage
\pagenumbering{arabic}

\section*{Course-note alignment}
\addcontentsline{toc}{section}{Course-note alignment}

This solution follows the course Notes directly:

\begin{itemize}
  \item Chapter 1: the model $\bm y=\bm X\bm\beta+\bm\varepsilon$,
  least squares, $\Rtwo$, and the partial $F$-test;
  \item Chapter 2: linearising an exponential model with $\log y$ and
  back-transforming to the original scale;
  \item Chapter 3: ridge penalties and the distinction between penalising and
  not penalising the intercept;
  \item Chapter 4: logistic probabilities, log-odds, odds ratios, and the
  threshold rule $\widehat\pi\ge 0.5$.
\end{itemize}

\section{Problem 1: Multiple linear regression}

For five students, let $x_1$ be self-study time, $x_2$ be tutoring time,
and $y$ be the test score. The assumed population model is
\[
  Y_i=\beta_0+\beta_1x_{i1}+\beta_2x_{i2}+\varepsilon_i,
  \qquad \varepsilon_i\overset{\mathrm{iid}}{\sim}N(0,\sigma^2).
\]

\subsection{Part (a): Least-squares estimates and R-squared}

\subsubsection*{Core idea}

Least squares finds the plane whose vertical squared residuals have the
smallest possible total. In matrix notation,
\[
  \bm y=\bm X\bm\beta+\bm\varepsilon,
  \qquad
  \widehat{\bm\beta}=(\bm X^{\mathsf T}\bm X)^{-1}
  \bm X^{\mathsf T}\bm y.
\]

The design matrix and response vector are
\[
\bm X=
\begin{bmatrix}
1&5&1\\
1&9&0\\
1&7&7\\
1&9&7\\
1&3&8
\end{bmatrix},
\qquad
\bm y=
\begin{bmatrix}
5\\6\\9.5\\10\\7.5
\end{bmatrix}.
\]
Therefore,
\[
\bm X^{\mathsf T}\bm X=
\begin{bmatrix}
5&33&23\\
33&245&141\\
23&141&163
\end{bmatrix},
\qquad
\bm X^{\mathsf T}\bm y=
\begin{bmatrix}
38\\258\\201.5
\end{bmatrix}.
\]

Solving the normal equations
$\bm X^{\mathsf T}\bm X\widehat{\bm\beta}
=\bm X^{\mathsf T}\bm y$ gives
\[
\boxed{
\widehat{\bm\beta}=
\begin{bmatrix}
\widehat\beta_0\\[2pt]\widehat\beta_1\\[2pt]\widehat\beta_2
\end{bmatrix}
=
\begin{bmatrix}
1.819205\\0.486520\\0.558644
\end{bmatrix}}
\]
and hence
\[
  \boxed{\widehat y=1.819205+0.486520x_1+0.558644x_2.}
\]

Holding tutoring time fixed, one more hour of self-study is associated with
an estimated $0.4865$-point increase in score. Holding self-study time fixed,
one more tutoring hour is associated with an estimated $0.5586$-point increase.

Using the course-note definitions,
\[
  \SSE=\sum_{i=1}^n(y_i-\widehat y_i)^2=0.281267,
  \qquad
  \SST=\sum_{i=1}^n(y_i-\overline y)^2=18.7.
\]
Thus
\[
  \boxed{
  \Rtwo=1-\frac{\SSE}{\SST}
  =1-\frac{0.281267}{18.7}
  =0.984959.}
\]
About $98.50\%$ of the observed score variation in this five-student sample is
explained by the two-predictor fitted model.

\subsubsection*{R verification}

\begin{lstlisting}
study <- data.frame(
  x1 = c(5, 9, 7, 9, 3),
  x2 = c(1, 0, 7, 7, 8),
  y  = c(5, 6, 9.5, 10, 7.5)
)

fit_full <- lm(y ~ x1 + x2, data = study)
X <- model.matrix(fit_full)
Y <- matrix(study$y, ncol = 1)

crossprod(X)                       # X'X
crossprod(X, Y)                    # X'y
solve(crossprod(X), crossprod(X,Y))

SSE <- sum(resid(fit_full)^2)
SST <- sum((study$y - mean(study$y))^2)
1 - SSE / SST
\end{lstlisting}

\subsection{Part (b): Compare the reduced and full models}

The reduced model is $Y_i=\beta_0+\beta_1x_{i1}+\varepsilon_i$.
The full model adds $x_2$. The hypotheses are
\[
  H_0:\beta_2=0
  \qquad\text{versus}\qquad
  H_1:\beta_2\ne 0.
\]

From Chapter 1, the partial $F$ statistic is
\[
F=
\frac{(\SSE_R-\SSE_F)/(df_{E,R}-df_{E,F})}
     {\SSE_F/df_{E,F}}.
\]
Substituting the supplied ANOVA quantities,
\[
F=
\frac{(16.7941-0.2813)/(3-2)}{0.2813/2}
=117.42,
\qquad p=0.008409.
\]

At $\alpha=0.05$, $p<\alpha$, so reject $H_0$. Tutoring time supplies
significant additional information after self-study time is already in the
model. The more appropriate of the two candidates is
\[
  \boxed{Y\sim x_1+x_2.}
\]

\begin{lstlisting}
fit_reduced <- lm(y ~ x1, data = study)
anova(fit_reduced, fit_full)
\end{lstlisting}

\subsection{Part (c): Prediction}

For $x_1=5$ and $x_2=6$,
\begin{align*}
\widehat y
&=1.819205+0.486520(5)+0.558644(6)\\
&=\boxed{7.603669}.
\end{align*}

\begin{lstlisting}
predict(fit_full, newdata = data.frame(x1 = 5, x2 = 6))
# 7.603669
\end{lstlisting}

\subsection*{Casio fx-580VN X check}

In \texttt{MENU $\to$ Matrix}, enter $\bm X$ as MatA and $\bm y$ as
MatB. Compute $\operatorname{Trn}(\text{MatA})\text{MatA}$ and
$\operatorname{Trn}(\text{MatA})\text{MatB}$, then evaluate

\[
(\bm X^{\mathsf T}\bm X)^{-1}\bm X^{\mathsf T}\bm y.
\]
The displayed
coefficient vector should be approximately

\[
(1.819205,0.486520,0.558644)^{\mathsf T}.
\]

\section{Problem 2: Exponential regression}

The proposed model is
\[
  Y_i=\exp(\beta_0+\beta_1x_i+\varepsilon_i),
  \qquad \varepsilon_i\sim N(0,\sigma^2).
\]

\subsection{Part (a): Linearisation and estimation}

Taking natural logarithms produces a model that is linear in its parameters:
\[
  Z_i=\log Y_i=\beta_0+\beta_1x_i+\varepsilon_i.
\]
The transformed observations are
\[
\begin{array}{c|cccc}
x_i&1&2&3&4\\ \hline
y_i&2.98&4.65&5.00&6.50\\
z_i=\log y_i&1.091923&1.536867&1.609438&1.871802
\end{array}
\]

For simple linear regression,
\[
  \widehat\beta_1=\frac{S_{xz}}{S_{xx}},
  \qquad
  \widehat\beta_0=\overline z-\widehat\beta_1\overline x,
\]
where
\[
  S_{xx}=\sum_{i=1}^n(x_i-\overline x)^2,
  \qquad
  S_{xz}=\sum_{i=1}^n(x_i-\overline x)(z_i-\overline z).
\]
The estimates are
\[
  \boxed{\widehat\beta_0=0.924456,
  \qquad \widehat\beta_1=0.241221.}
\]
Therefore,
\[
  \boxed{\log\widehat y=0.924456+0.241221x.}
\]
After back-transformation,
\[
  \widehat y
  =e^{0.924456}e^{0.241221x}
  =2.52050(1.27280)^x.
\]
Thus one additional year multiplies the fitted population by
$e^{0.241221}=1.27280$, an estimated increase of about $27.28\%$.

\begin{lstlisting}
population <- data.frame(
  year = 1:4,
  population = c(2.98, 4.65, 5.00, 6.50)
)
fit_exp <- lm(log(population) ~ year, data = population)
coef(fit_exp)
exp(coef(fit_exp)["year"])
\end{lstlisting}

\subsection{Part (b): Predict the fifth year}

\begin{align*}
\log\widehat y_5
&=0.924456+0.241221(5)=2.130559,\\
\widehat y_5
&=e^{2.130559}=\boxed{8.419576\text{ million people}.}
\end{align*}

This plug-in back-transformation estimates the conditional median on the
original scale. Under an exact lognormal model, the conditional mean would be
$\exp(\widehat\eta+\widehat\sigma^2/2)$; the standard exam answer is the
requested plug-in prediction above.

\begin{lstlisting}
exp(predict(fit_exp, newdata = data.frame(year = 5)))
# 8.419576
\end{lstlisting}

\subsection*{Casio fx-580VN X check}

Compute $\log y_i$ and use \texttt{MENU $\to$ Statistics $\to$ A+BX} with
$x_i$ and $\log y_i$. The calculator should return
$a\approx0.924456$ and $b\approx0.241221$. Evaluate
$\exp(a+5b)$ to obtain $8.419576$.

\section{Problem 3: Logistic regression}

Let
\[
  \pi(\bm x)=P(Y=1\mid\bm X=\bm x).
\]
Following Chapter 4, logistic regression models the log-odds:
\[
  \logit\{\pi(\bm x)\}
  =\log\left(\frac{\pi(\bm x)}{1-\pi(\bm x)}\right)
  =\beta_0+\beta_1x_1+\beta_2x_2+\beta_3x_3.
\]

\subsection{Part (a): Estimated model and interpretation}

The fitted model is
\[
\boxed{
\log\left(\frac{\widehat\pi}{1-\widehat\pi}\right)
=-3.449548+0.002294x_1+0.777014x_2-0.560031x_3.}
\]
Equivalently, if
\[
\widehat\eta=-3.449548+0.002294x_1+0.777014x_2-0.560031x_3,
\]
then
\[
  \widehat\pi=\frac{e^{\widehat\eta}}{1+e^{\widehat\eta}}
  =\frac{1}{1+e^{-\widehat\eta}}.
\]

Holding the other predictors fixed, a one-unit increase in $x_j$ multiplies
the odds of class 1 by $e^{\widehat\beta_j}$:
\[
\begin{array}{c|c|c|l}
\text{Predictor}&\widehat\beta_j&e^{\widehat\beta_j}&\text{Interpretation}\\ \hline
x_1&0.002294&1.002297&\text{odds increase by about }0.23\%\\
x_2&0.777014&2.174968&\text{odds are multiplied by about }2.175\\
x_3&-0.560031&0.571191&\text{odds decrease by about }42.9\%
\end{array}
\]

All three reported slope p-values are below $0.05$, so each predictor is
statistically significant at the $5\%$ level within this fitted model.

\subsection{Part (b): Predicted probabilities and classes}

Use the default threshold $\tau=0.5$:
\[
\widehat Y=
\begin{cases}
1,&\widehat\pi\ge0.5,\\
0,&\widehat\pi<0.5.
\end{cases}
\]

\begin{center}
\begin{tabular}{ccccc}
\toprule
\textbf{Row}&$\widehat\eta$&$\widehat\pi$&
$\widehat\pi\ge0.5$&\textbf{Class}\\
\midrule
1&-1.058484&0.257599&\text{No}&0\\
2&-0.940450&0.280809&\text{No}&0\\
3& 0.362391&0.589619&\text{Yes}&1\\
4&-0.684196&0.335326&\text{No}&0\\
\bottomrule
\end{tabular}
\end{center}

Hence the predicted classes are
\[
  \boxed{(0,0,1,0).}
\]

\begin{lstlisting}
b <- c(-3.449548, 0.002294, 0.777014, -0.560031)
new <- data.frame(
  x1 = c(640, 600, 700, 620),
  x2 = c(3.35, 3.62, 3.56, 3.17),
  x3 = c(3, 3, 1, 2)
)

eta <- drop(cbind(1, as.matrix(new)) %*% b)
p <- plogis(eta)
cbind(new, eta, p, class = as.integer(p >= 0.5))
\end{lstlisting}

\subsection*{Casio fx-580VN X check}

For each row, compute $\widehat\eta$ and then
$1/(1+e^{-\widehat\eta})$. For row 1 the display should show
$\widehat\eta\approx-1.05848$ and $\widehat\pi\approx0.25760$, hence class 0.

\section{Problem 4: Ridge estimator proof}

\subsection*{What must be shown}

For $\lambda>0$, minimise
\[
L(\beta_0,\beta_1)
=\sum_{k=1}^n(y_k-\beta_0-\beta_1x_k)^2
+\lambda(\beta_0^2+\beta_1^2).
\]
Unlike the practical \texttt{glmnet} model, this question explicitly penalises
both the intercept and slope.

Define the centered sums
\[
S_{xx}=\sum_{k=1}^n(x_k-\overline x)^2,
\qquad
S_{xy}=\sum_{k=1}^n(x_k-\overline x)(y_k-\overline y).
\]

\subsection*{Proof strategy}

Differentiate $L$ with respect to both parameters. Solve the first normal
equation for $\widehat\beta_0$, substitute it into the second, and rewrite the
uncentered sums using $S_{xx}$ and $S_{xy}$.

\begin{proof}
Differentiating with respect to $\beta_0$ gives
\begin{align*}
\frac{\partial L}{\partial\beta_0}
&=-2\sum_{k=1}^n(y_k-\beta_0-\beta_1x_k)+2\lambda\beta_0.
\end{align*}
At the minimiser, this derivative is zero, so
\begin{align*}
0
&=-\sum_{k=1}^n y_k+n\widehat\beta_0
  +\widehat\beta_1\sum_{k=1}^n x_k+\lambda\widehat\beta_0\\
&=-n\overline y+(n+\lambda)\widehat\beta_0
  +n\overline x\widehat\beta_1.
\end{align*}
Therefore,
\begin{equation}
\boxed{
\widehat\beta_0
=\frac{n}{n+\lambda}
  (\overline y-\widehat\beta_1\overline x).}
\label{eq:beta0}
\end{equation}

Next,
\begin{align*}
\frac{\partial L}{\partial\beta_1}
&=-2\sum_{k=1}^n x_k(y_k-\beta_0-\beta_1x_k)
  +2\lambda\beta_1.
\end{align*}
Setting this equal to zero yields
\begin{equation}
\sum_{k=1}^n x_ky_k
=n\overline x\widehat\beta_0
+\widehat\beta_1\left(\sum_{k=1}^n x_k^2+\lambda\right).
\label{eq:secondnormal}
\end{equation}
Substitute \eqref{eq:beta0} into \eqref{eq:secondnormal}:
\begin{align*}
\sum_{k=1}^n x_ky_k
&=n\overline x\frac{n}{n+\lambda}
  (\overline y-\widehat\beta_1\overline x)
  +\widehat\beta_1\left(\sum_{k=1}^n x_k^2+\lambda\right).
\end{align*}
Collecting the $\widehat\beta_1$ terms gives
\begin{align}
\widehat\beta_1
\left[
\sum_{k=1}^n x_k^2+\lambda
-\frac{n^2}{n+\lambda}\overline x^{,2}
\right]
&=
\sum_{k=1}^n x_ky_k
-\frac{n^2}{n+\lambda}\overline x\,\overline y.
\label{eq:uncentered}
\end{align}

Use the identities
\[
\sum_{k=1}^n x_ky_k=S_{xy}+n\overline x\,\overline y,
\qquad
\sum_{k=1}^n x_k^2=S_{xx}+n\overline x^{,2}.
\]
The right-hand side of \eqref{eq:uncentered} becomes
\begin{align*}
S_{xy}+n\overline x\,\overline y
-\frac{n^2}{n+\lambda}\overline x\,\overline y
&=S_{xy}
+\frac{n\lambda}{n+\lambda}\overline x\,\overline y.
\end{align*}
The bracket on the left-hand side becomes
\begin{align*}
S_{xx}+n\overline x^{,2}+\lambda
-\frac{n^2}{n+\lambda}\overline x^{,2}
&=S_{xx}+\lambda
+\frac{n\lambda}{n+\lambda}\overline x^{,2}\\
&=S_{xx}
+\lambda\left(1+\frac{n}{n+\lambda}\overline x^{,2}\right).
\end{align*}
Consequently,
\begin{equation}
\boxed{
\widehat\beta_1=
\frac{
S_{xy}+\dfrac{n\lambda}{n+\lambda}\overline x\,\overline y
}{
S_{xx}+\lambda\left(1+\dfrac{n}{n+\lambda}\overline x^{,2}\right)
}.}
\label{eq:beta1}
\end{equation}
Equations \eqref{eq:beta0} and \eqref{eq:beta1} are the required results.
\end{proof}

Because $L$ is a strictly convex quadratic for $\lambda>0$, its Hessian is
positive definite; therefore the stationary point above is the unique global
minimiser.

\subsection*{R verification of the proof}

The following code compares the derived formulas with the penalised matrix
solution. Here the $2\times2$ identity matrix penalises both coefficients.

\begin{lstlisting}
x <- c(1, 2, 4, 5)
y <- c(2, 3, 7, 8)
lambda <- 2
n <- length(y)
x_bar <- mean(x); y_bar <- mean(y)
Sxx <- sum((x - x_bar)^2)
Sxy <- sum((x - x_bar) * (y - y_bar))

b1 <- (Sxy + n*lambda*x_bar*y_bar/(n+lambda)) /
      (Sxx + lambda*(1 + n*x_bar^2/(n+lambda)))
b0 <- n*(y_bar - b1*x_bar)/(n+lambda)

X <- cbind(1, x)
matrix_answer <- solve(crossprod(X) + lambda*diag(2),
                       crossprod(X, y))
rbind(derived = c(b0, b1), matrix = matrix_answer)
#          [,1] [,2]
# derived  0.333  1.5
# matrix   0.333  1.5
\end{lstlisting}

\section*{Final answer summary}
\addcontentsline{toc}{section}{Final answer summary}

\begin{center}
\begin{tabular}{>{\raggedright\arraybackslash}p{0.22\textwidth}
                >{\raggedright\arraybackslash}p{0.68\textwidth}}
\toprule
\rowcolor{navy!10}\textbf{Item}&\textbf{Answer}\\
\midrule
1(a)&$\widehat y=1.819205+0.486520x_1+0.558644x_2$,
$\Rtwo=0.984959$\\
1(b)&Use the full model; $F=117.42$, $p=0.008409$\\
1(c)&$\widehat y(5,6)=7.603669$\\
2(a)&$\log\widehat y=0.924456+0.241221x$\\
2(b)&$\widehat y_5=8.419576$ million\\
3(a)&$\logit(\widehat\pi)=-3.449548+0.002294x_1+0.777014x_2-0.560031x_3$\\
3(b)&Predicted classes $(0,0,1,0)$\\
4&Equations \eqref{eq:beta0} and \eqref{eq:beta1}\\
\bottomrule
\end{tabular}
\end{center}

\end{document}

